\SourcePage{306}{309}
\section{Atomic Energy Levels}

In the nonrelativistic approximation, the stationary states of an atom are
determined by the Schr\"odinger equation for a system of electrons moving in
the Coulomb field of the nucleus and interacting electrically with one
another. Electron-spin operators do not occur in this equation at all. For a
system of particles in a centrally symmetric external field, both the total
orbital angular momentum $L$ and the parity of the state are conserved.
Accordingly, every stationary state of an atom has a definite value of $L$
and a definite parity. The coordinate-space wavefunctions of stationary
states of identical particles also possess a definite permutation symmetry.
Section~63 established that, for a system of electrons, every particular type
of permutation symmetry---that is, every particular Young diagram---is
associated with a particular value of the total spin of the system. The total
electron spin $S$ therefore provides a further label for each stationary
atomic state.

For fixed $S$ and $L$, an energy level is degenerate with respect to the
possible spatial directions of the vectors $\mathbf S$ and $\mathbf L$. The
degeneracy factors associated with their directions are $2S+1$ and $2L+1$,
respectively. Hence the complete degeneracy of a level having specified $L$
and $S$ is $(2L+1)(2S+1)$.

The electromagnetic interaction between electrons does, however, contain
relativistic effects that depend on their spins. Because of these effects, the
atomic energy depends not only on the magnitudes of $\mathbf L$ and $\mathbf
S$, but also on their relative orientation. Strictly speaking, once the
relativistic interactions are included, neither the orbital angular momentum
$L$ nor the spin $S$ of the atom is separately conserved. Only the combined
angular momentum
\[
\mathbf J=\mathbf L+\mathbf S
\]
remains conserved. This is a universal exact law that follows from the
isotropy of space for a closed system. Exact energy levels must consequently
be labelled by the value $J$ of the total angular momentum.

\SourcePage{307}{310}
When the relativistic effects are comparatively small, as is often the case,
they may instead be treated as a perturbation. Under this perturbation, the
degenerate level with given $L$ and $S$ splits into a number of distinct,
closely spaced levels having different values of the total angular momentum
$J$.

To first order, these levels are obtained from the corresponding secular
equation (Section~39). Their zeroth-order wavefunctions are certain linear
combinations of the wavefunctions belonging to the original degenerate level
with the specified values of $L$ and $S$.

Within this approximation, the magnitudes of the orbital angular momentum and
the spin may thus still be regarded as conserved, although their directions
may not. The levels can therefore continue to be labelled by $L$ and $S$ as
well.

Thus relativistic effects divide a level of specified $L$ and $S$ into levels
with several different values of $J$.
Such splitting is termed the level's \emph{fine structure} (or \emph{multiplet splitting}).

The
allowed values of $J$ run from $L+S$ down to $|L-S|$. A level with given $L$
and $S$ therefore yields $2S+1$ distinct levels when $L>S$, or $2L+1$ when
$L<S$. Each resulting level remains degenerate with respect to the directions
of $\mathbf J$, with degeneracy $2J+1$. As a check, summing $2J+1$ over every
allowed value of $J$ gives $(2L+1)(2S+1)$, as required.

Atomic energy levels, also called atomic \emph{spectral terms}, are customarily
denoted by symbols analogous to those used for individual-particle states of
definite angular momentum (Section~32). Different values of the total orbital
angular momentum $L$ are represented by capital Latin letters according to
the correspondence
\[
\begin{array}{r@{\;}cccccccccccc}
L={}&0&1&2&3&4&5&6&7&8&9&10&\ldots\\
&S&P&D&F&G&H&I&K&L&M&N&\ldots
\end{array}
\]
The number $2S+1$, called the \emph{multiplicity} of the term, is placed as a
superscript to the left of the symbol. One should remember, however, that this
number equals the number of fine-structure components only when
$L\geqslant S$\footnote{For $2S+1=1,2,3,\ldots$, the levels are called
singlet, doublet, triplet, and so on, respectively.}. The value of
\SourcePage{308}{311}
the total angular momentum $J$ is placed as a subscript on the right. Thus
${}^{2}P_{1/2}$ and ${}^{2}P_{3/2}$ denote levels with $L=1$, $S=1/2$, and
$J=1/2,3/2$.
